%
%============= CHAPTER BREAK =================%
%
\chapter{Claudius Ptolemy (ca 127--178 CE)}%
\vspace*{3mm}%
\justifying%
%
\noindent{}Claudius Ptolemy is estimated to have lived between 100 and 178 CE.\endnote{%
Robbins 1964:vi.} %
The astronomical observations in his \textit{Almagest} %
date between 127 and 151 CE, %
and its composition was mostly if not entirely post-138 CE. %
The \textit{Tetrabiblos} mentions the \textit{Almagest}, %
and thus the \textit{Tetrabiblos} dates sometime after 151 CE.\endnote{Robbins 1964:vi, viii.} \\[1mm]%

\noindent{}\textbf{ἐξορκιστής} / \textit{eksorkistēs} / “exorcist” may be a commonly known word today, %
but its first appearance in non-Christian Greek literature according to the TLG is %
in the astrologer Ptolemy. %

\vspace*{3mm} 
 \begin{quote}\textls[-15]{ἰδίως δὲ πάλιν, ἐὰν ἡ σελήνη τὸν πρακτικὸν τόπον ἐπισχῇ % 
 τὸν ἀπὸ συνόδου δρόμον ποιουμένη σὺν τῷ τοῦ Ἑρμοῦ %
 ... ποιεῖ ... ἐν δὲ Ζυγῷ καὶ Κριῷ καὶ Λέοντι θεολήπτους, ὀνειροκρίτας, \textbf{ἐξορκιστάς}.} (\textit{Anth}. 4.4.11) \\[1mm] %
 Characteristically again, if the moon holds the active place, %
 making course on a conjunction with Mercury %
 ... this produces ... in Libra and Aries and Leo them who are God-possessed, dream-readers, \textbf{exorcists}. %
 \end{quote} % 
 \vspace*{3mm}
 
 \noindent{}This evidently new Greek word also appears in Acts 19:13. %
 The passage opens by claiming God worked miracles through Paul (19:11), %
 healing diseases even by way of cloth that had touched him (19:12). %
 Then:

\vspace*{3mm} 
 \begin{quote}\textls[-15]{ἐπεχείρησαν δέ τινες καὶ
 τῶν περιερχομένων Ἰουδαίων \textbf{ἐξορκιστῶν} ὀνομάζειν ἐπὶ
 τοὺς ἔχοντας τὰ πνεύματα τὰ πονηρὰ τὸ ὄνομα τοῦ
 κυρίου Ἰησοῦ λέγοντες, Ὁρκίζω ὑμᾶς τὸν Ἰησοῦν ὃν
 Παῦλος κηρύσσει.} \clearpage %
 Some of the itinerant Judean \textbf{exorcists} attempted to name the name of the Lord Jesus %
 over those who had the evil spirits, saying, %
 I oath-bind you by Jesus, whom Paul preaches. \end{quote} % 
 \vspace*{3mm}

\noindent{}This immediately leads to the exorcists (seven sons of Sceva) %
getting pummeled and stripped naked by the demon-possessed man (19:14). %

An epigram falsely ascribed to Lucian of Samosata and likely dated to after the 2nd century CE %
also uses this term, in quite grotesque fashion: %

\vspace*{3mm} 
 \begin{quote}Δαίμονα πολλὰ λαλῶν ὀζόστομος \textbf{ἐξορκιστὴς} \newline
   ἐξέβαλ’, οὐχ ὅρκων, ἀλλὰ κόπρων δυνάμει. (\textit{Epigram}. 11.427) \\[1mm] %
 Speaking to many demons, a foul-breathed \textbf{exorcist} \newline
 cast them out, not with oaths, but instead with the power of shits. \end{quote} % 
 \vspace*{3mm} 

\noindent{}The cumulative evidence points to Acts making use of a lexical convention %
of the mid- to late-2nd century CE at the earliest. \\[1.5mm] %

\noindent{}\textbf{μεθοδεία} / \textit{methodeia} / “method” %
first occurs in a chronologically secure setting %
in Ptolemyʼs \textit{Syntaxis} or \textit{Almagest}, %
completed by 151 CE.\endnote{%
1.2:581.} %

\vspace*{3mm} 
 \begin{quote}τὰ δὲ πέμπτα σελίδια γέγονεν ἡμῖν ὑπὲρ τοῦ καὶ
 τὰς ἐν ταῖς ἄλλαις τῶν ἐκκέντρων παρόδοις συνισταμένας %
 κατὰ πλάτος παραχωρήσεις διευκρινεῖν ἐκ τῆς τῶν
 παρατιθεμένων ἑξηκοστῶν \textbf{μεθοδείας}. \\[1mm] %
 Now the fifth little columns were made by us concerning %
 the occurrences in the other passages of cardinal points %
 according to width, to clarify extensions \textbf{by method} %
 of the sixty presented alongside. \end{quote} % 
 \vspace*{3mm}
%

\noindent{}It also appears in the [post-]Aurelian \textit{Partitiones} of [Ps-]Herodian: %

\vspace*{3mm} 
 \begin{quote}\textls[-10]{Τὰ ἀπὸ τῶν διὰ τοῦ ευω ῥημάτων, διὰ τοῦ εια γινόμενα θηλυκὰ, %
 διὰ τῆς ει διφθόγγου γράφονται· οἷον· 
 λατρεύω, λατρεία· μεσιτεύω, μεσιτεία· πρεσβεύω, πρεσβεία· %
 καρυκεύω, καρυκεία· βλακεύω, βλακεία· μεθοδεύω, \textbf{μεθοδεία}· %
 καὶ τὰ ὅμοια ὡσαύτως. (Boissonade 1819:221)} \\[1mm] %
 The words from those formed with the ευω, %
 being feminine with the εια, %
 are written with the ει diphthong, such as: %
 to worship, worship; to mediate, mediation; to represent, representation; %
 to sauce up, sauce; to be lazy, laziness; to do methodically, \textbf{method}; %
 and the like similarly. %
 \end{quote} % 
 \vspace*{3mm}

Among early Christian works, %
this word first appears in a parallel portion of Marcionite Laodiceans and %
canonical Ephesians, with only modest [bracketed] %
differences between them (Lao 6:12 // Eph 6:12):  

\vspace*{3mm} 
 \begin{quote} ἐνδύσασθε τὴν πανοπλίαν τοῦ θεοῦ %
 [πρὸς τὸ δύνασθαι ὑμᾶς] τὸ στῆναι πρὸς τὰς \textbf{μεθοδείας} τοῦ διαβόλου· \\[1mm] %
 Be clothed in the panoply of the God [in order for you to be able] to stand %
 up to the \textbf{methods} of the devil. 
 \end{quote} % 
 \vspace*{3mm}

\noindent{}In a verse possibly missing from Marcionite Laodiceans, %
the same term appeared earlier in canonical Ephesians:  

\vspace*{3mm} 
 \begin{quote}ἵνα μηκέτι ὦμεν νήπιοι, %
 κλυδωνιζόμενοι καὶ περιφερόμενοι παντὶ ἀνέμῳ %
 τῆς διδασκαλίας ἐν τῇ κυβείᾳ τῶν ἀνθρώπων ἐν πανουργίᾳ %
 πρὸς τὴν \textbf{μεθοδείαν} τῆς πλάνης. \\[1.5mm] %
 so that we may no longer be infants, %
 wave-tossed and carried round by every wind %
 of the teaching in the trickery of the humans %
 in craftiness unto the \textbf{method} of the deception. \end{quote} % 
 \vspace*{3mm}

 One explicit and one supplied instance of this lemma turned up in searches of digitized Greek inscriptions. %
 Milet VI,2 565 uses μεθοδείᾳ explicitly.\endnote{%
Milet VI,2 565 — PH351291 — TM947458 — Ionia, Miletus, Kalabaktepe — limestone sarcophagus epitaph.} %
 Though a precise date is not provided in the PHI or TM databases, %
 the text itself mentions Fl[avius] Antiochianus, %
 perhaps one of the two famous persons by that name who lived in the 3rd century CE.\endnote{%
 \textit{PLRE} 1:70.} %
 The supplied hit (IGLSyr 4 1593) is clearly Christian and dates between 383--395 CE.\endnote{%
 IGLSyr 4 1593 — PH243843 — Apamene; see \url{epigraphy.packhum.org/text/243843}.} %

 Another neologism for which Ptolemy is arguably the earliest attestor is %
 κοσμοκράτωρ / \textit{kosmokratōr} / “world-ruler”, %
 but in our view Antigonus provides the earliest attestation, %
 and so this lemma is explored in the chapter devoted to Antigonus. %
 Here in the conclusion of the chapter on Ptolemy, %
 we note that two new and/or highly distinctive terms (κοσμοκράτωρ, μεθοδεία) of Ptolemy %
 are found adjacent to each other in Marcionite Laodiceans and canonical Ephesians (6:11--12). %
 This double-finding based on computational queries of all Greek lemmas inside the NT %
 within a broader diachronic view strongly favors the overshadowing influence of the famous astrologer. %
 Certainly, the responses in Marcionite Laodiceans and canonical Ephesians strike an adversarial tone, %
 but appropriation and critique often go hand in hand. %
 Among known candidates for authorship of Marcionite Laodiceans, Valentinus is as good as any. %
 Even if it is presently difficult to be certain of this, %
 we have good reason to be confident about the rhetorical opponent in view: the astrologer Ptolemy. %